\documentclass[11pt,a4paper]{article}

\usepackage[utf8]{inputenc}
\usepackage[T1]{fontenc}
\usepackage[margin=1in]{geometry}
\usepackage{amsmath,amssymb}
\usepackage{booktabs}
\usepackage{graphicx}
\usepackage{xcolor}
\usepackage{hyperref}
\usepackage{fancyhdr}
\usepackage{titlesec}
\usepackage{enumitem}
\usepackage{caption}
\usepackage{array}
\definecolor{accentviolet}{HTML}{6D54C9}
\definecolor{accentblue}{HTML}{2E6FB8}
\definecolor{good}{HTML}{1F9D68}
\definecolor{textdim}{HTML}{5B5480}

\hypersetup{
  colorlinks=true,
  linkcolor=accentviolet,
  citecolor=accentviolet,
  urlcolor=accentblue,
  pdftitle={Kshanik --- Technical Product Report},
  pdfauthor={Team Transience},
  pdfsubject={Physics-Informed Yield Surrogate for a Non-Isothermal Plug-Flow Reactor}
}

\titleformat{\section}{\normalfont\Large\bfseries\color{accentviolet}}{\thesection}{1em}{}
\titleformat{\subsection}{\normalfont\large\bfseries}{\thesubsection}{1em}{}

\pagestyle{fancy}
\fancyhf{}
\renewcommand{\headrulewidth}{0.4pt}
\fancyhead[L]{\small\textit{Kshanik --- Technical Product Report}}
\fancyhead[R]{\small Team Transience}
\fancyfoot[C]{\thepage}

\newcommand{\rmse}{\text{RMSE}}

\title{
  {\Huge \textbf{Kshanik}}\\[0.15em]
  {\normalsize \textit{(kshanik --- Hindi for ``transient, momentary'')}}\\[0.5em]
  {\Large Physics-Informed Yield Surrogate for a Non-Isothermal Plug-Flow Reactor}\\[0.6em]
  {\large Technical Product Report}
}
\author{
  Team Transience\\[0.3em]
  \small Vaishnav Aditya \quad $\cdot$ \quad Preenithi Varshana \quad $\cdot$ \quad Kritika Pandey\\[0.3em]
  \small Fugacity 2026 ML Hackathon, IIT Kharagpur
}
\date{\small \today}

\begin{document}

\maketitle

\begin{abstract}
\noindent
We present Kshanik, a physics-informed machine learning surrogate for predicting the yield of a
non-isothermal, continuous-flow plug-flow reactor (PFR) operating a series reaction network
$A \rightarrow B \rightarrow C$. Rather than fitting a black-box regressor directly to five raw
operating conditions, we first calibrate a mechanistic ODE model --- coupled mole balances,
Arrhenius kinetics, and a jacket energy balance --- to the 150 officially provided training
observations via multi-start nonlinear least squares. This mechanistic model's prediction then
serves as a physically grounded prior mean for a Gaussian Process, trained in a logit-warped
target space to correctly respect the bounded, zero-inflated nature of the yield variable. This
small-data model achieves a 5$\times$10 repeated cross-validated $\rmse$ of \textbf{8.79} on the
150-row dataset. Recognizing that a hackathon-scale dataset is not representative of a production
deployment, we additionally built a second, scale-oriented model: a two-stage CatBoost
classifier--regressor trained on the real data augmented with 5{,}000 synthetic observations
generated by numerically integrating our own fitted physics model across the full operating
envelope, achieving $\rmse = 11.35$ and designed to improve as real production data accumulates.
Both models, all code, and this report are available in the accompanying GitHub repository.
\end{abstract}

\vspace{1em}
\noindent\textbf{Project links.}\\
GitHub repository: \url{https://github.com/adityakalagatoori/TRANSIENCE-IITK}\\
Project brief (GitHub Pages): \url{https://adityakalagatoori.github.io/TRANSIENCE-IITK/}

\section{Introduction and Problem Statement}

The Fugacity 2026 ML Hackathon (IIT Kharagpur) problem statement asks participants to build a
data-driven surrogate model that predicts the \texttt{overall\_yield} of a continuous-flow,
non-isothermal reactor, given five operating conditions:

\begin{itemize}[nosep]
  \item \texttt{flow\_rate\_L\_min} --- volumetric feed flow rate (L/min)
  \item \texttt{concentration\_mol\_L} --- inlet concentration of reactant $A$ (mol/L)
  \item \texttt{inlet\_temperature\_K} --- feed temperature entering the reactor (K)
  \item \texttt{length\_m} --- reactor length (m)
  \item \texttt{jacket\_temperature\_K} --- external heating/cooling jacket temperature (K)
\end{itemize}

\noindent The underlying chemistry is a series-parallel reaction network:
\begin{equation}
  A \xrightarrow{k_1} B \xrightarrow{k_2} C
\end{equation}
where $B$ is the desired product and $C$ is an undesired side product formed by further
conversion of $B$. The evaluation metric is root-mean-squared error (RMSE) on 50 held-out test
rows, scored blind against a hidden physics-based ground truth.

\subsection{Data Provenance}

All real data used in this project --- the 150-row \texttt{train\_dataset.csv} and the 50-row
\texttt{test\_dataset.csv} --- was provided directly by the official hackathon problem statement
and dataset release. No external or scraped datasets were used for training or evaluation. The
only additional data used anywhere in this project is a 5{,}000-row \emph{synthetic} dataset,
which we generated ourselves (Section~\ref{sec:catboost}) by numerically simulating our own
fitted mechanistic model across the operating envelope --- it is not sourced from any external
dataset and exists solely to test whether our modeling approach scales beyond the 150-row
hackathon regime.

\subsection{Why This Problem Is Hard}

Exploratory analysis of the 150 training rows reveals two properties that make direct black-box
regression difficult:

\begin{enumerate}[nosep]
  \item \textbf{Zero-inflation.} 37 of 150 rows (24.7\%) have \texttt{overall\_yield} exactly equal
  to zero, and 57 of 150 (38\%) have yield below 1\%. This is not measurement noise --- it reflects
  a genuine physical regime in which the side reaction $B \rightarrow C$ fully consumes the
  desired product before the reactor exit.
  \item \textbf{Weak raw-feature correlation.} No single raw input correlates strongly with yield;
  the strongest single-variable Pearson correlation (jacket temperature) is $|r| = 0.64$. Table~\ref{tab:corr}
  summarizes all five.
\end{enumerate}

\begin{table}[h]
\centering
\caption{Pearson correlation of each raw feature with \texttt{overall\_yield} (train set, $n=150$).}
\label{tab:corr}
\begin{tabular}{lc}
\toprule
\textbf{Feature} & \textbf{Correlation with yield} \\
\midrule
\texttt{flow\_rate\_L\_min}       & $+0.04$ \\
\texttt{concentration\_mol\_L}    & $+0.01$ \\
\texttt{inlet\_temperature\_K}    & $-0.41$ \\
\texttt{length\_m}                & $+0.08$ \\
\texttt{jacket\_temperature\_K}   & $-0.50$ \\
\bottomrule
\end{tabular}
\end{table}

\noindent This indicates that yield is governed by a non-linear combination of the raw inputs
(a Damk\"{o}hler-number-like quantity), not by any one variable directly ---
motivating a physics-first rather than purely data-driven approach.

\section{Mechanistic Foundation: The Governing Physical Model}
\label{sec:mechanistic}

Instead of treating the five inputs as opaque regressors, we model the reactor's governing
physics directly and calibrate its unknown parameters to the training data.

\subsection{Mole Balances}

For a plug-flow reactor with residence-time coordinate $\tau$, the steady-state mole balances for
the series reaction are:
\begin{align}
  \frac{dC_A}{d\tau} &= -k_1(T)\, C_A \\
  \frac{dC_B}{d\tau} &= k_1(T)\, C_A - k_2(T)\, C_B
\end{align}
with $C_A(0) = C_{A,0}$ (the inlet concentration feature) and $C_B(0) = 0$.

\subsection{Arrhenius Kinetics}

Both rate constants follow the Arrhenius law:
\begin{equation}
  k_i(T) = k_{0,i} \exp\!\left(-\frac{E_{a,i}}{RT}\right), \qquad i \in \{1, 2\}
\end{equation}
where $R = 8.314\ \text{J/(mol\,K)}$.

\subsection{Non-Isothermal Energy Balance}

Since \texttt{inlet\_temperature\_K} and \texttt{jacket\_temperature\_K} are independently
specified and frequently differ substantially in the dataset, the reactor is explicitly
non-isothermal. We follow the standard non-isothermal PFR energy balance
(Fogler, \emph{Elements of Chemical Reaction Engineering}, Ch.~8), which couples jacket heat
exchange with the heat generated or consumed by each reaction step:
\begin{equation}
  \frac{dT}{d\tau} = \frac{U_a (T_j - T) + (-\Delta H_1)\, r_1 + (-\Delta H_2)\, r_2}{\rho C_p}
\end{equation}
where $r_1 = k_1(T) C_A$ and $r_2 = k_2(T) C_B$ are the two reaction rates, $U_a$ is a lumped
heat-transfer coefficient, $\Delta H_1, \Delta H_2$ are the reaction enthalpies, and $\rho C_p$
is a lumped heat-capacity term.

An earlier version of this model included only the jacket heat-exchange term (equivalent to
assuming both reactions are athermal). We identified this as a simplification relative to the
standard textbook energy balance and added the explicit reaction-heat terms, which measurably
improved the mechanistic fit (Section~\ref{sec:mech-results}).

\subsection{Residence-Time Proxy}

Reactor cross-sectional area is not provided in the dataset, so we treat
$\tau_{\max} = \texttt{length\_m} / \texttt{flow\_rate\_L\_min}$ as a residence-time \emph{proxy},
with the unknown proportionality constant (cross-sectional area) absorbed into the fitted
pre-exponential factors $k_{0,1}, k_{0,2}$.

\subsection{Parameter Estimation}

Eight effective parameters --- $k_{0,1}, E_{a,1}, k_{0,2}, E_{a,2}, U_a, \Delta H_1, \Delta H_2,
\rho C_p$ --- are recovered from the 150 training rows via multi-start nonlinear least squares
(\texttt{scipy.optimize.least\_squares}, trust-region-reflective method), minimizing the sum of
squared errors between the simulated final yield $100 \cdot C_B(\tau_{\max}) / C_{A,0}$ and the
observed \texttt{overall\_yield} across all training rows. The ODE system is integrated with a
fixed-step 4th-order Runge--Kutta scheme (50 steps), chosen after diagnosing that adaptive solvers
(LSODA, Radau, RK23) could stall indefinitely at certain parameter-space corners explored during
optimization, where reaction rates and thermal derivatives become numerically extreme. The
fixed-step scheme, combined with physically-motivated caps on temperature, reaction rate, and
temperature-derivative magnitude, guarantees bounded computation time for any parameter draw
within the optimizer's search bounds.

Multi-start optimization (6--10 random initializations per fit) was used throughout to reduce the
risk of convergence to a poor local optimum, since the $(k_0, E_a)$ parameter landscape is
non-convex.

\subsection{Mechanistic Fit Results}
\label{sec:mech-results}

The fitted mechanistic model achieves a training RMSE of \textbf{12.17} with no machine learning
component whatsoever. Critically, the fitted parameters satisfy $E_{a,2} > E_{a,1}$: the side
reaction's activation energy exceeds the desired reaction's. This is not an assumption --- it is
recovered purely from the 150 training rows --- and it directly explains the zero-yield cliff
observed in the data: as temperature rises, the side reaction rate constant $k_2$ grows faster
than $k_1$, so at sufficiently high temperature the side reaction dominates and consumes
essentially all of product $B$ before the reactor exit.

\section{Small-Data Model: Warped Gaussian Process Regression}
\label{sec:gpr}

With only 150 labeled rows, we selected Gaussian Process Regression (GPR) as the primary
statistical model, since GPR is well suited to the small-data regime and provides calibrated
predictive uncertainty.

\subsection{Physics-Engineered Features}

Rather than feeding the GP the five raw operating conditions directly, we derive a compact set of
physically motivated features from the fitted mechanistic model:

\begin{itemize}[nosep]
  \item $\text{Da}_1 = k_1(T_{\text{avg}})\, \tau$, $\text{Da}_2 = k_2(T_{\text{avg}})\, \tau$
  --- Damk\"{o}hler numbers for each reaction step, and their log-transforms
  \item Selectivity ratio $k_2/k_1$ (log-transformed)
  \item \texttt{concentration\_mol\_L} (raw)
  \item $\Delta T = T_{\text{jacket}} - T_{\text{inlet}}$
\end{itemize}

This five-dimensional feature set collapses the original five raw inputs into physically
interpretable quantities that directly correspond to the drivers of conversion and selectivity in
series-reaction kinetics.

\subsection{Warped-Target Formulation}

The target variable \texttt{overall\_yield} is bounded on $[0, 100]$ and has substantial
probability mass at the boundary (Section~\ref{sec:mechanistic}). A standard GP with an additive
Gaussian-noise likelihood is statistically inappropriate for such a target: it can (and empirically
does) predict outside $[0, 100]$, and it treats the boundary pile-up as ordinary noise rather than
a structural feature of the response. This is precisely the setting addressed by the
censored-regression (Tobit) literature and, in a tractable non-parametric form, by
\emph{warped Gaussian processes} \cite{snelson2003warped}.

We apply a logit warp to the normalized target:
\begin{equation}
  z = \text{logit}\!\left(\frac{y}{100}\right) = \ln\!\left(\frac{y/100}{1 - y/100}\right),
  \qquad y \in [\epsilon, 100-\epsilon]
\end{equation}
with a small clipping margin $\epsilon = 0.3$ to keep $z$ finite. The mechanistic model's own
prediction $\mu(x)$ is passed through the same warp, $\mu_z(x) = \text{logit}(\mu(x)/100)$, and a
Mat\'{e}rn-5/2-kernel GP is trained on the residual $z - \mu_z(x)$ in this transformed space.
Predictions are inverted back via the logistic (sigmoid) function:
\begin{equation}
  \hat{y} = 100 \cdot \frac{1}{1 + e^{-\hat{z}}}
\end{equation}
which automatically respects the $[0,100]$ boundary for any GP output, by construction.

\subsection{Model Comparison}

We evaluated three variants of the GPR pipeline via repeated 5-fold cross-validation (10 repeats,
50 folds total), summarized in Table~\ref{tab:gpr-variants}.

\begin{table}[h]
\centering
\caption{Cross-validated RMSE for GPR pipeline variants (150 real rows, 5$\times$10 repeated K-fold).}
\label{tab:gpr-variants}
\begin{tabular}{lcc}
\toprule
\textbf{Variant} & \textbf{Mean RMSE} & \textbf{Std. Dev.} \\
\midrule
Mechanistic model alone (no ML) & 12.17 & --- \\
Single-stage GPR (additive residual) & 10.14 & 4.14 \\
Two-stage (classifier + GPR) & 10.85 & 5.06 \\
\textbf{Warped-logit GPR (selected)} & \textbf{8.79} & \textbf{3.87} \\
\bottomrule
\end{tabular}
\end{table}

\noindent Two findings are worth noting. First, the warped-logit formulation outperforms the
plain additive-residual GP by a meaningful margin (8.79 vs.\ 10.14), consistent with the
hypothesis that respecting the bounded/censored likelihood matters more than additional model
complexity. Second, we tested a two-stage architecture --- a binary classifier to separate the
``collapsed'' ($y \approx 0$) regime before applying the GP to the ``active'' regime --- and found
it performed \emph{worse} than the single-stage warped model, both in mean RMSE and in
cross-fold variance. With only 150 rows, the classifier introduces its own misclassification
error near the collapse boundary, and the mechanistic mean function already captures the
collapse physically (via $\text{Da}_2 \to$ large), making the extra stage redundant rather than
corrective. This is a deliberate negative result we report for completeness: it directly informed
our final architecture choice.

The warped-logit GPR is our official small-data submission, achieving $\rmse = 8.79$.

\section{Large-Data Model: CatBoost with Synthetic Augmentation}
\label{sec:catboost}

A 150-row training set is representative of a hackathon constraint, not of a production plant
deployment. To demonstrate that our approach scales beyond the officially provided dataset, we
built a second model explicitly designed for a larger-data regime.

\subsection{Synthetic Data Generation}

Using our fitted mechanistic model (Section~\ref{sec:mechanistic}) as a physics engine, we
generated 5{,}000 additional synthetic operating conditions via Latin Hypercube Sampling across
the five input ranges (with a 5\% margin beyond the observed train/test envelope), then simulated
each through the fitted ODE to obtain a physically-consistent synthetic yield. This synthetic set
exhibits a zero-yield fraction of 21.8--40\%, close to the 24.7\% observed in the real data,
indicating that the mechanistic model reproduces the real collapse-regime balance reasonably
well when extrapolated across the operating envelope.

We emphasize that this synthetic dataset is not an external data source: it is generated entirely
in-house from our own calibrated physics model, specifically to test scalability, and it is
never used to validate the final reported metrics (see below).

\subsection{Two-Stage Architecture}

The large-data model uses the same physics-engineered feature set as the GPR model (extended with
the raw inputs and additional Damk\"{o}hler-derived quantities, 13 features total), and a
two-stage CatBoost architecture:

\begin{enumerate}[nosep]
  \item A CatBoost \emph{classifier} predicts whether a row falls in the collapsed ($y < 0.5\%$)
  or active yield regime, trained on real rows (weight 1.0) plus synthetic rows (weight 0.15).
  \item A CatBoost \emph{regressor}, trained the same way but restricted to rows in the active
  regime, predicts the continuous yield for rows the classifier calls active. Rows called
  collapsed are predicted as zero.
\end{enumerate}

Synthetic rows are deliberately down-weighted (0.15$\times$) relative to real rows so the model
anchors primarily on ground truth, using the synthetic data mainly to regularize the yield surface
between sparsely sampled real observations --- particularly around the sharp collapse boundary.

\subsection{Validation Methodology}

All cross-validation is performed on real data only: synthetic rows are added to each fold's
\emph{training} split but never enter a held-out \emph{test} fold, since their ``ground truth''
originates from the same mechanistic model being evaluated --- including them in validation would
be circular. Table~\ref{tab:catboost-variants} summarizes results.

\begin{table}[h]
\centering
\caption{Cross-validated RMSE for CatBoost pipeline variants (evaluated on real data only, 5$\times$10 repeated K-fold).}
\label{tab:catboost-variants}
\begin{tabular}{lcc}
\toprule
\textbf{Variant} & \textbf{Mean RMSE} & \textbf{Std. Dev.} \\
\midrule
Real data only (no augmentation), single-stage & 15.24 & 3.17 \\
Real + synthetic, single-stage & 12.67 & 3.35 \\
\textbf{Real + synthetic, two-stage (selected)} & \textbf{11.35} & \textbf{3.65} \\
\bottomrule
\end{tabular}
\end{table}

\noindent Synthetic augmentation improves RMSE by roughly 17\% over real-data-only training
(15.24 $\to$ 12.67), and the two-stage classifier-gated architecture provides a further
improvement to 11.35. This is the inverse of the finding in Section~\ref{sec:gpr}: with only 150
real rows, a classifier stage hurt the GPR model, but with the additional synthetic data available
to the CatBoost pipeline, the same architectural idea helps. We interpret this as evidence that
the two-stage architecture is fundamentally sound, but data-hungry --- exactly the profile expected
of a model designed to improve further as real production data accumulates.

\section{Results Summary}

Table~\ref{tab:summary} consolidates all reported results.

\begin{table}[h]
\centering
\caption{Summary of all model variants and their role.}
\label{tab:summary}
\begin{tabular}{p{4.2cm}p{3.7cm}cc}
\toprule
\textbf{Model} & \textbf{Role} & \textbf{RMSE} & \textbf{Data used} \\
\midrule
Mechanistic ODE (physics only) & Foundation for both models & 12.17 & 150 real rows \\
Warped-logit GPR & \textbf{Official small-data submission} & \textbf{8.79} & 150 real rows \\
CatBoost, two-stage + synthetic & Scale-oriented model & 11.35 & 150 real + 5{,}000 synthetic \\
\bottomrule
\end{tabular}
\end{table}

Approximating $\rmse$ as an $R^2$-equivalent accuracy relative to the variance of the training
target ($\sigma_y^2 = 38.3^2$), the warped-logit GPR corresponds to $\approx 94.8\%$ of target
variance explained, and the CatBoost scale-model to $\approx 91.3\%$.

\section{Discussion: Two Models, One Physics Foundation}

A central design decision in this project was to build \emph{two} models rather than one, each
targeting a different data regime, but both grounded in the same calibrated mechanistic ODE:

\begin{itemize}[nosep]
  \item The \textbf{warped-logit GPR} is tuned specifically for the exact 150-row dataset provided
  by the hackathon. It is our official Phase~1 submission.
  \item The \textbf{CatBoost two-stage model} is engineered for scale: it is trained on an
  order-of-magnitude larger, physics-consistent synthetic dataset, uses a production-friendly
  gradient-boosted architecture (fast inference, straightforward retraining pipeline), and is
  designed to improve as real operating data accumulates from an actual plant deployment ---
  something a Gaussian Process, whose training cost scales cubically with dataset size, cannot do
  as gracefully.
\end{itemize}

Both models are reproducible from the training data alone: the mechanistic parameters, the
physics-engineered features, and the final predictions are all generated by scripts in the
accompanying repository, with no manually tuned or hard-coded values beyond model
hyperparameters.

\section{Reproducibility and Code Availability}

All code, both trained models' final predictions, the fitted mechanistic parameters, the
generated synthetic dataset, and this report are available at:

\begin{center}
\url{https://github.com/adityakalagatoori/TRANSIENCE-IITK}
\end{center}

\noindent The repository README documents the full file layout and exact commands to reproduce
every result in this report from the original \texttt{train\_dataset.csv} and
\texttt{test\_dataset.csv} files.

\begin{thebibliography}{9}

\bibitem{fogler}
H.~S. Fogler. \emph{Elements of Chemical Reaction Engineering}, 4th/5th ed., Chapter~8:
Steady-State Non-Isothermal Reactor Design. Prentice Hall.

\bibitem{snelson2003warped}
E.~Snelson, C.~E. Rasmussen, Z.~Ghahramani. ``Warped Gaussian Processes.''
\emph{Advances in Neural Information Processing Systems (NeurIPS)}, 2003.

\bibitem{tobin1958}
J.~Tobin. ``Estimation of Relationships for Limited Dependent Variables.''
\emph{Econometrica}, 26(1), 24--36, 1958.

\end{thebibliography}

\end{document}
